\begin{textsample}{POS Dim 2 – 1950 – Score 13.00 – 1950/tv\_com\_1950\_97.txt}  \label{ex:f2_pos_086}
In black and white, a middle-aged man in work clothes and a soft brimmed cap sits in profile as if driving.
He is hunched forward, squeezed close to a small \textbf{steering} wheel, with his knees drawn high and his shoulders rounded.
His face looks strained and uncomfortable.
The background streaks \textbf{past} in a blur of houses, \textbf{trees}, and roadway, making him appear to be moving, though only the driver ’s seat area and \textbf{steering} wheel are clearly visible.
Gradually, a second image of the same man overlaps him : this version sits farther back and more upright, with both hands relaxed on the wheel.
As the uncomfortable figure fades away, the man now appears cheerful and at ease, smiling as he \textbf{drives} against the blurred \textbf{outdoor} background.
The scene shifts into an indoor showroom or studio display.
A light-colored Ford pickup \textbf{truck} is parked on a glossy floor, angled so its front and driver ’s side are visible.
A well-dressed male presenter, middle-aged, with dark hair, a dark suit, white shirt, patterned tie, and pocket square, stands beside the driver ’s door.
Behind the \textbf{truck} are wall posters and tall glass-block-style windows ; one poster includes Ford \textbf{truck} imagery.
The presenter gestures toward the cab and the door, then moves closer to the handle.
He presses or uses the small horizontal door handle, and the driver ’s door swings open.
The opening is shown \textbf{broad} and tall, with the presenter standing partly behind the door to emphasize its size.
A close view focuses on the door hinge area : a thick, curved metal hinge arm connects the door to the cab pillar, with large bolts visible above and below.
The hinge is shown as sturdy and mechanical, with the open door held away from the cab.
Back at the \textbf{truck}, the presenter stands in front of the open doorway, the door swung wide toward the front fender.
He steps toward the cab and climbs in, using the wide opening.
He bends at the waist, places a hand inside for \textbf{balance}, and slides onto the bench seat.
Once seated, he faces outward through the open doorway, one hand resting on the seat cushion, the \textbf{steering} wheel visible to his right inside the cab.
The focus moves closer to the seat.
The presenter ’s hand rests on a \textbf{broad}, dark, rounded bench cushion.
He presses and pats the cushion surface, demonstrating its thickness and softness.
The image transitions into a cutaway illustration of the seat : part of the upholstery is removed to reveal a grid of sinuous metal springs in the seat base and backrest, with a rounded tubular frame along the edges.
Large white arrows point to the spring areas, and bold on-screen text reads “ NON-SAG SPRINGS. ” Another close diagram of the seat base shows the spring pattern and a small central mechanism ; the words “ SEAT SHOCK SNUBBER ” appear in large white letters, with an arrow pointing to the mechanism.
A simple side-profile graphic of a \textbf{truck} seat then appears on a plain background.
The backrest is shown in several overlapping positions, with a curved arrow above it to indicate it can tilt.
Bold text below reads “ ADJUSTABLE BACKREST. ” The cutaway seat returns, now showing the spring structure beside the upholstered cushion.
Bold text across the top reads “ FOAM RUBBER CUSHIONS, ” with a white arrow pointing toward the cushion area.
The diagram fades back into the real cab, where the presenter ’s hand again presses the dark bench cushion.
The cushion slightly changes shape under his hand, visually reinforcing softness.
He then shifts his posture, drawing one knee up slightly and turning his upper body toward the dashboard.
His right hand gestures across the interior toward the windshield and instrument panel.
The \textbf{steering} wheel, dashboard, windshield frame, and open door frame surround him, giving a clear sense of the cab ’s spacious interior.
The presenter leans forward and points toward the large front windshield from inside the cab.
A front view of the \textbf{truck} follows, with the presenter visible through the \textbf{broad}, one-piece windshield.
The hood fills the lower part of the view, and the windshield spans nearly the full width of the cab.
A rectangular rear-view mirror hangs near the upper center.
The driver ’s door remains open at the left side of the \textbf{truck}.
The presenter sits behind the wheel and gestures while looking outward through the glass.
A rear view of the cab shows the large rectangular rear window.
The presenter is visible inside, turned slightly as if looking back through it.
The rear window takes up much of the back wall of the cab and is divided by a narrow vertical frame, emphasizing the amount of glass.
The view returns to the open driver ’s doorway, where the presenter sits comfortably, smiling and speaking toward the outside.
A large Ford \textbf{truck} poster hangs on the wall behind him.
He gestures with one hand, then reaches toward the door and pulls it inward.
The door swings across the opening ; he remains seated inside, leaning an arm on the lowered side window opening.
The side window is \textbf{broad} and deep, with the presenter ’s face framed by the glass and door.
He looks out through the side opening, then glances downward toward the controls and back up again, maintaining a relaxed seated posture.
The image dissolves from the showroom cab to a Ford \textbf{truck} on an open road.
For a moment, the presenter ’s face and cab window overlap with the \textbf{outdoor} scene, then the road scene becomes clear.
A large Ford \textbf{truck} with a rounded 1950s cab and a long, high-sided cargo body \textbf{drives} along a \textbf{two-lane} \textbf{rural} road.
The roadway has painted center lines ; utility poles, bare \textbf{trees}, fields, and low roadside \textbf{brush} stretch into the distance.
The \textbf{truck} moves steadily toward and \textbf{past} the viewer ’s position, its grille, headlights, rounded hood, and heavy-duty wheels visible.
The cargo body behind the cab is tall and rectangular, with vertical side panels.
A closer front view shows three men seated side by side across the cab, visible through the wide windshield.
The driver sits on the right side of the image with both hands near the \textbf{steering} wheel.
Two passengers sit beside him, one in the middle and one on the left.
All three appear adult men, wearing work-style caps or dark headwear, and they sit upright with enough shoulder room to be clearly separated.
The wide windshield frames them together, visually demonstrating the cab ’s width.
A car \textbf{passes} in the adjacent lane on the right side of the road in the background, while the Ford \textbf{truck} continues forward.
The scene returns to the wider road view, with the \textbf{truck} moving along the \textbf{rural} \textbf{highway}.
Large bold white lettering appears over the moving \textbf{truck} and road : “ FORD ” at the top, “ ECONOMY ” in a slanted script-style word beneath it, “ \textbf{TRUCKS} ” in large block letters, and “ for ’ 53 ” below.
The \textbf{truck} remains visible behind the text as it continues down the road.
The background darkens gradually while the lettering stays centered, and the commercial fades out with the words “ FORD ECONOMY \textbf{TRUCKS} for ’ 53 ” still visible.

% matched lemmas: balance, broad, brush, drive, highway, outdoor, pass, past, rural, steering, tree, truck, two-lane
\end{textsample}
